\documentclass[a4paper,11pt]{report}
\usepackage[dutch]{babel}
\usepackage{geometry}
\usepackage{parskip}
\usepackage{iflang}

\begin{document}

%%=============================================================================
%% Methodologie
%%=============================================================================

\chapter{\IfLanguageName{dutch}{Methodologie}{Methodology}}%
\label{ch:methodologie}

In dit hoofdstuk wordt de algemene aanpak van het onderzoek toegelicht. Om een antwoord te formuleren op de onderzoeksvraag \textit{``Welke machine learning modellen zijn het meest geschikt voor het detecteren van misinformatie op Mastodon?''} is gekozen voor een kwantitatieve, experimentele onderzoeksmethode. Het onderzoek is opgedeeld in vier opeenvolgende fasen, die hieronder worden beschreven.

\section{Fase 1: Literatuurstudie}
De eerste fase bestond uit een grondige literatuurstudie om de theoretische basis van het onderzoek te leggen.
\begin{itemize}
    \item \textbf{Doelstelling:} Inzicht verwerven in de technische architectuur van Mastodon (ActivityPub-protocol), de huidige state-of-the-art in automatische detectie van misinformatie, en de werking van de geselecteerde algoritmen (SVM, Logistic Regression en BERT).
    \item \textbf{Methode:} Desk research via wetenschappelijke databanken zoals Google Scholar, IEEE Xplore en arXiv. Er is specifiek gezocht naar recente publicaties over content moderatie in het Fediverse.
    \item \textbf{Deliverable:} De bevindingen uit deze fase vormen de basis voor Hoofdstuk 2 (Stand van zaken).
\end{itemize}

\section{Fase 2: Dataverzameling en Pre-processing}
Omdat er geen kant-en-klare dataset bestaat voor misinformatie op Mastodon, stond het samenstellen van een geschikte dataset centraal in deze fase.
\begin{itemize}
    \item \textbf{Doelstelling:} Het verkrijgen van voldoende gelabelde data om machine learning modellen te trainen en te testen in een Mastodon-context.
    \item \textbf{Methode:}
    \begin{enumerate}
        \item \textbf{Trainingdata:} Gebruik van de LIAR-dataset, een benchmark voor fake news detectie. De labels zijn vereenvoudigd naar een binaire classificatie (waar/onwaar) om compatibiliteit te garanderen.
        \item \textbf{Testdata (Mastodon):} Ontwikkeling van een Python-script dat via de publieke Mastodon API berichten verzamelt. Er zijn minimaal 500 berichten verzameld, specifiek gezocht op hashtags gerelateerd aan actualiteit.
        \item \textbf{Verificatie:} De Mastodon-berichten zijn gelabeld aan de hand van betrouwbare factcheck-bronnen zoals \textit{VRT NWS Factcheck} en \textit{PolitiFact}.
    \end{enumerate}
    \item \textbf{Deliverable:} Een gecombineerde, opgeschoonde dataset klaar voor modeltraining.
\end{itemize}

\section{Fase 3: Modelontwikkeling}
In deze fase zijn de daadwerkelijke detectiemodellen gebouwd en getraind.
\begin{itemize}
    \item \textbf{Doelstelling:} Het implementeren van drie vergelijkbare classificatiemodellen.
    \item \textbf{Methode:}
    \begin{itemize}
        \item \textbf{Traditionele modellen:} Implementatie van Support Vector Machines (SVM) en Logistic Regression met behulp van de \textit{scikit-learn} library. Tekst is omgezet naar vectoren met TF-IDF.
        \item \textbf{Deep Learning model:} Fine-tuning van een pre-trained BERT-model (Bidirectional Encoder Representations from Transformers) met behulp van de \textit{Hugging Face} library. Dit model kan context en woordvolgorde beter begrijpen dan de traditionele methoden.
    \end{itemize}
    \item \textbf{Deliverable:} Drie getrainde en functionele modellen.
\end{itemize}

\section{Fase 4: Evaluatie en Analyse}
De laatste fase richtte zich op het beantwoorden van de deelvragen door de resultaten van de modellen te analyseren.
\begin{itemize}
    \item \textbf{Doelstelling:} Bepalen welk model het beste presteert en welke factoren hier invloed op hebben.
    \item \textbf{Methode:}
    \begin{itemize}
        \item \textbf{Kwantitatieve evaluatie:} De modellen zijn vergeleken op basis van standaard metrieken: F1-score, Precision en Recall.
        \item \textbf{Factoranalyse:} Er is specifiek gekeken naar de invloed van berichtlengte (kort vs. lang), de aanwezigheid van bronvermeldingen en het taalgebruik op de classificatiefouten.
    \end{itemize}
    \item \textbf{Deliverable:} De resultaten en conclusies zoals beschreven in Hoofdstuk 4 en 5 van deze scriptie.
\end{itemize}

\end{document}